\textbf{Phần 1. Trắc nghiệm lựa chọn}
\setcounterpageref{ex}{0}

\begin{ex} %[4D12N63]
Muốn nhân hai phân số, ta
\choice
{Quy đồng mẫu số rồi nhân tử số}
{\True Lấy tử số nhân với tử số, mẫu số nhân với mẫu số}
{Lấy tử số nhân mẫu số chéo}
{Giữ nguyên mẫu số, nhân các tử số}
\loigiai{
Lấy tử số nhân với tử số, mẫu số nhân với mẫu số.\\
\textbf{Đáp án của câu hỏi là B}
}
\end{ex}

\begin{ex} %[4D12N63]
Kết quả của phép tính $\dfrac{2}{3} \times \dfrac{4}{5}$ là
\choice
{$\dfrac{6}{8}$}
{$\dfrac{8}{8}$}
{\True $\dfrac{8}{15}$}
{$\dfrac{10}{12}$}
\loigiai{
$\dfrac{2 \times 4}{3 \times 5} = \dfrac{8}{15}$.\\
\textbf{Đáp án của câu hỏi là C}
}
\end{ex}

\begin{ex} %[4D12H63]
Kết quả của phép tính $\dfrac{3}{8} \times 4$ là
\choice
{$\dfrac{12}{32}$}
{$\dfrac{7}{8}$}
{$3$}
{\True $\dfrac{3}{2}$}
\loigiai{
$\dfrac{3 \times 4}{8} = \dfrac{12}{8} = \dfrac{3}{2}$.\\
\textbf{Đáp án của câu hỏi là D}
}
\end{ex}

\begin{ex} %[4D12H63]
Kết quả của phép tính $\dfrac{5}{7} \times \dfrac{7}{5}$ là
\choice
{\True $1$}
{$\dfrac{35}{10}$}
{$0$}
{$\dfrac{25}{49}$}
\loigiai{
$\dfrac{5 \times 7}{7 \times 5} = 1$.\\
\textbf{Đáp án của câu hỏi là A}
}
\end{ex}

\begin{ex} %[4D12H63]
Một hình vuông có cạnh dài $\dfrac{3}{4}\text{ m}$. Diện tích hình vuông đó là
\choice
{$3\text{ m}^2$}
{$\dfrac{6}{8}\text{ m}^2$}
{\True $\dfrac{9}{16}\text{ m}^2$}
{$\dfrac{3}{16}\text{ m}^2$}
\loigiai{
$\dfrac{3}{4} \times \dfrac{3}{4} = \dfrac{9}{16}\text{ m}^2$.\\
\textbf{Đáp án của câu hỏi là C}
}
\end{ex}

\begin{ex} %[4D12H63]
Tính thuận tiện: $\dfrac{2}{5} \times \dfrac{3}{7} \times \dfrac{5}{2}$.
\choice
{$1$}
{\True $\dfrac{3}{7}$}
{$\dfrac{6}{35}$}
{$\dfrac{15}{14}$}
\loigiai{
$\left(\dfrac{2}{5} \times \dfrac{5}{2}\right) \times \dfrac{3}{7} = 1 \times \dfrac{3}{7} = \dfrac{3}{7}$.\\
\textbf{Đáp án của câu hỏi là B}
}
\end{ex}

\begin{ex} %[4D12V63]
Một hình chữ nhật có chiều dài $\dfrac{4}{5}\text{ m}$, chiều rộng $\dfrac{3}{8}\text{ m}$. Diện tích hình chữ nhật đó là
\choice
{$\dfrac{7}{13}\text{ m}^2$}
{$\dfrac{12}{40}\text{ m}^2$}
{$\dfrac{3}{5}\text{ m}^2$}
{\True $\dfrac{3}{10}\text{ m}^2$}
\loigiai{
$\dfrac{4}{5} \times \dfrac{3}{8} = \dfrac{12}{40} = \dfrac{3}{10}\text{ m}^2$.\\
\textbf{Đáp án của câu hỏi là D}
}
\end{ex}

\begin{ex} %[4D12V63]
Một tấm kính hình thoi có độ dài hai đường chéo là $\dfrac{3}{4}\text{ m}$ và $\dfrac{2}{5}\text{ m}$. Diện tích tấm kính đó là
\choice
{$\dfrac{3}{10}\text{ m}^2$}
{\True $\dfrac{3}{20}\text{ m}^2$}
{$\dfrac{6}{20}\text{ m}^2$}
{$\dfrac{3}{40}\text{ m}^2$}
\loigiai{
$S = \dfrac{1}{2} \times \dfrac{3}{4} \times \dfrac{2}{5} = \dfrac{3}{20}\text{ m}^2$.\\
\textbf{Đáp án của câu hỏi là B}
}
\end{ex}

\begin{ex} %[4D12V63]
Tính giá trị biểu thức: $\dfrac{3}{4} \times \dfrac{5}{6} + \dfrac{3}{4} \times \dfrac{1}{6}$.
\choice
{\True $\dfrac{3}{4}$}
{$\dfrac{1}{2}$}
{$1$}
{$\dfrac{5}{8}$}
\loigiai{
$\dfrac{3}{4} \times \left(\dfrac{5}{6} + \dfrac{1}{6}\right) = \dfrac{3}{4} \times 1 = \dfrac{3}{4}$.\\
\textbf{Đáp án của câu hỏi là A}
}
\end{ex}

\begin{ex} %[4D12V63]
Một ô tô mỗi phút đi được $\dfrac{4}{5}\text{ km}$. Hỏi trong $15$ phút ô tô đó đi được bao nhiêu ki-lô-mét?
\choice
{$10\text{ km}$}
{$14\text{ km}$}
{$15\text{ km}$}
{\True $12\text{ km}$}
\loigiai{
$\dfrac{4}{5} \times 15 = 12\text{ km}$.\\
\textbf{Đáp án của câu hỏi là D}
}
\end{ex}

\begin{ex} %[4D12V63]
Tính: $\dfrac{2 \times 3 \times 4}{3 \times 4 \times 5}$.
\choice
{$\dfrac{1}{5}$}
{$\dfrac{3}{5}$}
{\True $\dfrac{2}{5}$}
{$\dfrac{4}{5}$}
\loigiai{
Rút gọn cho $3 \times 4$ ta được $\dfrac{2}{5}$.\\
\textbf{Đáp án của câu hỏi là C}
}
\end{ex}

\begin{ex} %[4D12C63]
Tính tích: $P = \left(1 - \dfrac{1}{2}\right) \times \left(1 - \dfrac{1}{3}\right) \times \left(1 - \dfrac{1}{4}\right) \times \dots \times \left(1 - \dfrac{1}{10}\right)$.
\choice
{\True $\dfrac{1}{10}$}
{$\dfrac{1}{5}$}
{$\dfrac{1}{9}$}
{$\dfrac{2}{10}$}
\loigiai{
$P = \dfrac{1}{2} \times \dfrac{2}{3} \times \dfrac{3}{4} \times \dots \times \dfrac{9}{10} = \dfrac{1}{10}$.\\
\textbf{Đáp án của câu hỏi là A}
}
\end{ex}

\textbf{Phần 2. Trắc nghiệm đúng sai}
\setcounterpageref{ex}{0}

\begin{ex} %[4D12H63]
Xét tính đúng sai của các phép nhân phân số:
\choiceTF
{\True $\dfrac{1}{3} \times \dfrac{2}{5} = \dfrac{2}{15}$.}
{\True $\dfrac{3}{4} \times 0 = 0$.}
{$\dfrac{5}{7} \times 1 = \dfrac{7}{5}$.}
{\True $\dfrac{5}{7} \times 1 = \dfrac{5}{7}$.}
\loigiai{
\begin{itemize}
    \item $1/3 \times 2/5 = 2/15$. Đúng.
    \item Nhân với 0 bằng 0: Đúng.
    \item Nhân với 1 bằng chính nó $\neq 7/5$. Ý 3 Sai, Ý 4 Đúng.
\end{itemize}
\begin{itemize}
    \item a) \textbf{Đúng}.
    \item b) \textbf{Đúng}.
    \item c) \textbf{Sai}.
    \item d) \textbf{Đúng}.
\end{itemize}
}
\end{ex}

\begin{ex} %[4D12V63]
Một mảnh bìa hình chữ nhật có chiều dài $\dfrac{3}{4}\text{ m}$, chiều rộng $\dfrac{2}{3}\text{ m}$. Xét tính đúng sai:
\choiceTF
{\True Chu vi mảnh bìa là $\dfrac{17}{6}\text{ m}$.}
{\True Diện tích mảnh bìa là $\dfrac{1}{2}\text{ m}^2$.}
{Diện tích mảnh bìa là $\dfrac{6}{7}\text{ m}^2$.}
{\True Nếu tăng chiều dài gấp đôi thì diện tích tăng gấp đôi ($1\text{ m}^2$).}
\loigiai{
\begin{itemize}
    \item Chu vi: $(3/4 + 2/3) \times 2 = 17/12 \times 2 = 17/6\text{ m}$. Đúng.
    \item Diện tích: $3/4 \times 2/3 = 6/12 = 1/2\text{ m}^2$. Đúng.
    \item $1/2 \neq 6/7$. Sai.
    \item Gấp đôi dài thì diện tích gấp đôi $= 1\text{ m}^2$: Đúng.
\end{itemize}
\begin{itemize}
    \item a) \textbf{Đúng}.
    \item b) \textbf{Đúng}.
    \item c) \textbf{Sai}.
    \item d) \textbf{Đúng}.
\end{itemize}
}
\end{ex}

\textbf{Phần 3. Trả lời ngắn}
\setcounterpageref{ex}{0}

\begin{ex} %[4D12H63]
Tính tích: $\dfrac{3}{5} \times \dfrac{5}{9}$ (rút gọn tối giản dạng a/b).
\shortans{1/3}
\loigiai{
\textbf{Đáp số:} $1/3$\\
$15/45 = 1/3$.
}
\end{ex}

\begin{ex} %[4D12H63]
Tính tích: $\dfrac{7}{8} \times 16$.
\shortans{14}
\loigiai{
\textbf{Đáp số:} $14$\\
$7 \times 2 = 14$.
}
\end{ex}

\begin{ex} %[4D12V63]
Một hình vuông có cạnh dài $\dfrac{2}{5}\text{ m}$. Tính chu vi hình vuông đó (dạng a/b).
\shortans{8/5}
\loigiai{
\textbf{Đáp số:} $8/5$\\
$2/5 \times 4 = 8/5\text{ m}$.
}
\end{ex}

\begin{ex} %[4D12V63]
Tính diện tích hình chữ nhật có kích thước $\dfrac{5}{6}\text{ m}$ và $\dfrac{3}{10}\text{ m}$ (rút gọn tối giản dạng a/b).
\shortans{1/4}
\loigiai{
\textbf{Đáp số:} $1/4$\\
$15/60 = 1/4\text{ m}^2$.
}
\end{ex}

\begin{ex} %[4D12V63]
Tính thuận tiện: $\dfrac{4}{9} \times \dfrac{7}{11} + \dfrac{4}{9} \times \dfrac{4}{11}$ (dạng a/b).
\shortans{4/9}
\loigiai{
\textbf{Đáp số:} $4/9$\\
$4/9 \times 1 = 4/9$.
}
\end{ex}

\begin{ex} %[4D12C63]
Tính tích: $\dfrac{1}{2} \times \dfrac{2}{3} \times \dfrac{3}{4} \times \dfrac{4}{5} \times \dfrac{5}{6}$ (dạng a/b).
\shortans{1/6}
\loigiai{
\textbf{Đáp số:} $1/6$\\
$1/6$.
}
\end{ex}
